Theorem 4. Let $I, P, Q, R$, and $S$ be as in the preceding theorems, and let $A_{1}$ and $A_{2}$ be disjoint arcs in $\bar{I}$, such that $A_{1} \cap \operatorname{Fr} I=\{P\}$ and $A_{2} \cap \operatorname{Fr} I= \{R\}$. Then $S$ and $Q$ are in the frontier of the same component of $I-\left(A_{1} \cup A_{2}\right)$.

\begin{center}
\fbox{\parbox{0.82\linewidth}{\centering \emph{Figure 4.5 omitted in this rendered excerpt; see the raw book source above.}}}
\end{center}


Proof. In Figure 4.5, we have shown $\bar{I}$ as a rectangular region, with $P$ and $R$ as the midpoints of a pair of opposite sides. See Theorem 3.5 and Problem 3.5.

By a brick-decomposition of the plane we mean a collection $G=\left\{g_{i}\right\}$ of polyhedral disks ( $=2$-cells) such that (1) $\cup_{i=1}^{\infty} g_{i}=\mathbf{R}^{2}$, (2) if two sets $g_{i}$ and $g_{j}$ intersect, then their intersection is a broken line lying in the frontier of each of them, and (3) every point has a neighborhood which intersects at most three of the sets $g_{i}$. One way to get such a collection is to cut up the plane by horizontal lines and vertical segments so as to get a sort of "infinite brick wall," as shown in Figure 4.6.

\begin{center}
\fbox{\parbox{0.82\linewidth}{\centering \emph{Figure 4.6 omitted in this rendered excerpt; see the raw book source above.}}}
\end{center}


In general, given a set $M$ in a metric space $[X, d]$, the diameter $\delta M$ of $M$ is the least upper bound of the numbers $d(P, Q)(P, Q \in M)$. (Thus $\delta M$\\
may be $\infty$.) If $G$ is a collection of subsets of $X$, then the mesh of $G$ is the least upper bound of the numbers $\delta g(g \in G)$.

Evidently we can construct a brick-decomposition $G$ of $\mathbf{R}^{2}$ with mesh as small as we please. In any case, the union of any subcollection of $G$ is a 2 -manifold with boundary. In the present case, we use bricks which are rectangular regions, with sides parallel to the edges of $\operatorname{Fr} I$, such that $\bar{I}$ is the union of a subcollection of them. And we use bricks of sufficiently small diameter so that no one of them intersects both $A_{1}$ and $A_{2}$. (See Problem 1.12.)

Let $N$ be the union of all bricks in the decomposition that intersect $A_{1}$. Then $N$ is a 2-manifold with boundary, and so also is the set

$$
N^{\prime}=N \cap \bar{I} .
$$

Let $J$ be the component of $\operatorname{Fr} N^{\prime}$ that contains $P$. Then $J$ is a 1 -sphere. Let $B_{1}$ be the component of $J \cap \mathrm{Fr} I$ that contains $P$. Then $B_{1}$ is a broken line between two points $T$ and $U$, where $T, U \in \mathrm{Fr} I, T$ lies below $P$ and $R$, and $U$ lies above $P$ and $R$. Let $B_{2}$ be the other broken line between $T$ and $U$ in $J$. Let $V$ be the last point of $B_{2}$ (in the order from $T$ ) that lies in Fr $I$ and lies below $P$ and $R$; and let $W$ be the first point of $B_{2}$ that follows $V$ and lies in $\operatorname{Fr} I$. Then $W$ lies above $P$ and $R$ in $\operatorname{Fr} I$.

\begin{center}
\fbox{\parbox{0.82\linewidth}{\centering \emph{Figure 4.7 omitted in this rendered excerpt; see the raw book source above.}}}
\end{center}


Let $B$ be the broken line between $V$ and $W$ in $B_{2}$. Then $B \cap \operatorname{Fr} I= \{V, W\}$. Thus $V$ and $W$ lie in the boundary of the same component of $I-\left(A_{1} \cup A_{2}\right)$. Therefore $Q$ and $S$ have the same property. \(\square\)
